% 그림: 사람이 챙기지 않아도 도는 자동 장치
\documentclass[tikz,border=10pt]{standalone}
% 그림 공통 설정 (INSTALL.tex와 같은 글꼴·색)
\usepackage{kotex}
\setmainfont{Noto Sans CJK KR}[UprightFont={* DemiLight},BoldFont={* Bold}]
\setsansfont{Noto Sans CJK KR}[UprightFont={* DemiLight},BoldFont={* Bold}]
\setmonofont{D2Coding ligature}[Scale=0.92]
\setmainhangulfont{Noto Sans CJK KR}[UprightFont={* DemiLight},BoldFont={* Bold}]
\setsanshangulfont{Noto Sans CJK KR}[UprightFont={* DemiLight},BoldFont={* Bold}]
\setmonohangulfont{D2Coding ligature}[Scale=0.92]
\usetikzlibrary{positioning,fit,arrows.meta,backgrounds,calc,shapes.geometric}
\definecolor{accent}{HTML}{0E7C86}
\definecolor{ink}{HTML}{1F2328}
\definecolor{muted}{HTML}{6B7280}
\definecolor{msgbg}{HTML}{E6F3F4}
\definecolor{okc}{HTML}{2E9E5B}
\definecolor{warnc}{HTML}{CC3B3B}
\tikzset{
  box/.style={draw=#1!55!black,fill=#1!8,rounded corners=3pt,align=center,inner sep=5pt,text=ink},
  box/.default=gray,
  lab/.style={font=\footnotesize,text=muted,align=center},
  arr/.style={-{Latex[length=2.4mm]},thick,#1},
  arr/.default=ink,
  num/.style={circle,fill=accent,text=white,font=\scriptsize\bfseries,inner sep=1.2pt,minimum size=4mm},
}

\begin{document}
\begin{tikzpicture}[font=\small]
  \tikzset{trig/.style={box=gray,minimum width=34mm,minimum height=12mm},
           act/.style={box=violet,text width=56mm,minimum height=12mm},
           res/.style={box=green,minimum width=34mm,minimum height=12mm}}
  % 1 아침 브리핑
  \node[trig] (t1) at (0,0) {\textbf{매일 09:00}\\[-1pt]{\footnotesize 또는 PI 봇의 \texttt{/brief}}};
  \node[act]  (a1) at (58mm,0) {\textbf{아침 브리핑}\\[-1pt]{\footnotesize 스크립트가 kanban 보드·git 상태를 모으고 Claude가 정리 (PI 모델은 쓰지 않음)}};
  \node[res]  (r1) at (116mm,0) {\#hermes 또는\\보낸 채널에 게시};
  % 2 차단 감시
  \node[trig] (t2) at (0,-22mm) {\textbf{10분마다}\\[-1pt]{\footnotesize 같은 kanban 카드가 2번 막힘}};
  \node[act]  (a2) at (58mm,-22mm) {\textbf{차단 감시}\\[-1pt]{\footnotesize Claude가 원인 진단, 허용 범위 안의 해결은 바로 (실행 권한, 명세 불일치, kanban 카드 정리)}};
  \node[res]  (r2) at (116mm,-22mm) {\#hermes에 결과\\{\footnotesize ``연구자 결정 필요''만 답함}};
  % 3 PI 인계
  \node[trig] (t3) at (0,-44mm) {\textbf{PI 대화 압축}\\[-1pt]{\footnotesize 대화가 길어질 때마다}};
  \node[act]  (a3) at (58mm,-44mm) {\textbf{PI 대화 인계}\\[-1pt]{\footnotesize Claude가 인계 노트를 다시 씀. 세 번째 압축이면 다음 메시지에서 새 대화를 열고 노트를 넣어 줌}};
  \node[res]  (r3) at (116mm,-44mm) {PI가 규칙과\\진행 상황을 유지};
  \foreach \i in {1,2,3} {
    \draw[arr=muted] (t\i) -- (a\i);
    \draw[arr=muted] (a\i) -- (r\i);
  }
  \node[lab,font=\footnotesize\bfseries] at (0,12mm) {언제};
  \node[lab,font=\footnotesize\bfseries] at (58mm,12mm) {무엇을};
  \node[lab,font=\footnotesize\bfseries] at (116mm,12mm) {결과};
  \node[lab,anchor=north] at (58mm,-54mm) {설치: \texttt{extras/install\_extras.sh} \quad 매 턴 PI 규칙 리마인더(\texttt{souls/pi-rules.md})도 같은 플러그인이 붙입니다};
\end{tikzpicture}
\end{document}
